\documentclass[11pt]{article}
\usepackage[utf8]{inputenc}
\usepackage[T1]{fontenc}
\usepackage{amsmath,amssymb}
\usepackage{graphicx}
\usepackage{booktabs}
\usepackage{hyperref}
\usepackage[margin=1in]{geometry}
\usepackage{natbib}
\usepackage{authblk}

\title{Integrating Single-Cell RNA and TCR/BCR Sequencing with Mass Cytometry Reveals Dynamic Trajectories of Human Peripheral Immune Cells from Birth to Old Age}

\author[1]{Author A}
\author[1]{Author B}
\author[1,2]{Author C}
\affil[1]{Shanghai Institute of Immunology, Shanghai Jiao Tong University School of Medicine, Shanghai, China}
\affil[2]{Department of Immunology, Pudong New Area People's Hospital, Shanghai, China}

\date{}

\begin{document}
\maketitle

\begin{abstract}
Aging profoundly reshapes the human immune system, yet a comprehensive single-cell multi-omic atlas spanning the full human lifespan has been lacking. Here, we profiled peripheral immune cells from 220 healthy volunteers across 13 age groups (0 to $>$90 years) from the Shanghai Pudong Cohort using single-cell RNA sequencing coupled with T cell receptor (TCR) and B cell receptor (BCR) sequencing, mass cytometry (CyTOF), bulk RNA-seq, and flow cytometry. We identified dynamic changes in immune cell composition, with CD4$^+$ T cells declining from 38.2\% in neonates to 25.5\% in the oldest group, while CD8$^+$ T cells and NK cells expanded progressively. T cells were the most strongly affected lineage, with GNLY$^+$CD8$^+$ effector memory T cells exhibiting the highest clonal expansion (mean clonality 0.52 in the elderly versus 0.08 in children). We discovered a population of cytotoxic B cells enriched in children (5.2\% at 0--1 years versus 0.5\% at $>$80 years) expressing high levels of granzyme B. Finally, we constructed a single-cell immune aging clock using 156 features across 28 cell types that predicts chronological age with a Pearson correlation of 0.92 and mean absolute error of 5.8 years. This multi-omic atlas provides a comprehensive resource for understanding human immune aging.
\end{abstract}

\section{Introduction}

The immune system undergoes continuous remodeling throughout the human lifespan, a process broadly termed immunosenescence \citep{nikolich2018immune}. These changes contribute to the increased susceptibility of older adults to infections, reduced vaccine efficacy, and elevated risk of autoimmune and inflammatory diseases \citep{goronzy2013understanding}. While the broad outlines of immune aging are well established---including thymic involution, contraction of the naive T cell compartment, and expansion of memory and effector populations---the precise dynamics of these changes at single-cell resolution across the full human lifespan have remained poorly characterized.

Previous studies of immune aging have typically focused on limited age ranges, often comparing only young adults to elderly individuals, or have relied on bulk approaches that cannot capture the heterogeneity within immune cell populations \citep{alpert2019multi}. The advent of single-cell RNA sequencing (scRNA-seq) has enabled high-resolution profiling of immune cell states, but comprehensive lifespan studies integrating transcriptomics with immune repertoire sequencing remain scarce \citep{zheng2020landscape}. Furthermore, the dynamics of T cell and B cell receptor repertoires across the lifespan are poorly characterized at single-cell resolution, and it is unknown which immune cell subsets undergo the most intensive rewiring at specific life stages.

Mass cytometry (CyTOF) provides complementary high-dimensional protein-level information that can validate and extend transcriptomic findings \citep{bendall2011single}. Integrating these modalities offers a more complete picture of immune system dynamics than any single approach alone.

Here, we present a multi-omic single-cell atlas of peripheral immune cells spanning the full human lifespan, from birth to over 90 years of age. By profiling 220 healthy volunteers from the Shanghai Pudong Cohort across 13 age groups using five complementary technologies, we reveal dynamic trajectories of immune cell composition, transcriptomic states, clonal expansion, and cell--cell interactions. We identify GNLY$^+$CD8$^+$ effector memory T cells as the subset with the highest clonal expansion in aging, discover a previously unappreciated population of cytotoxic B cells enriched in children, and construct a single-cell immune aging clock that can assess individual immune status across life stages.

\section{Related Work}

Single-cell approaches have been increasingly applied to study immune aging. \citet{zheng2020landscape} performed scRNA-seq on peripheral blood mononuclear cells (PBMCs) from individuals of different ages but examined a limited age range. \citet{alpert2019multi} used mass cytometry to profile immune aging but lacked single-cell transcriptomic and repertoire data. Studies of T cell repertoire diversity have shown age-associated contraction \citep{britanova2014age}, but these relied on bulk TCR sequencing that cannot link clonotype to phenotype at the single-cell level.

The concept of an immune clock has been explored using bulk omics data. \citet{sayed2021immuneclock} developed an inflammatory aging clock from plasma proteomics, and \citet{alpert2019multi} proposed CyTOF-based aging signatures. However, no study has combined single-cell transcriptomics, repertoire sequencing, and protein-level measurements to build a comprehensive immune aging clock.

The discovery of unconventional B cell subsets has expanded our understanding of B cell biology \citep{glass2020unconventional}. Age-associated B cells accumulate in the elderly and are linked to autoimmunity, but the existence of cytotoxic B cell populations in early life has not been previously reported.

\section{Methods}

\subsection{Study cohort and sample collection}
We recruited 220 healthy volunteers from the Shanghai Pudong Cohort spanning 13 age groups: 0--1 year, 1--5 years, 6--12 years, 13--18 years, 19--30 years, 31--45 years, 46--60 years, 61--70 years, 71--80 years, and $>$80 years. Peripheral blood was collected and processed for multiple assays.

\subsection{Single-cell RNA sequencing with TCR/BCR profiling}
PBMCs were isolated by density gradient centrifugation and subjected to 10x Genomics Chromium Single Cell 5' v2 library preparation, enabling simultaneous capture of gene expression, TCR, and BCR sequences from the same cells. Libraries were sequenced on Illumina NovaSeq 6000 platforms.

\subsection{Mass cytometry (CyTOF)}
A panel of $>$40 metal-conjugated antibodies targeting major immune lineage and functional markers was used to profile PBMCs by CyTOF. Data were normalized using bead standards and analyzed using standard gating and unsupervised clustering approaches.

\subsection{Bulk RNA sequencing and flow cytometry}
Bulk RNA-seq was performed on sorted immune cell populations for validation. Flow cytometry with conventional fluorescent antibodies was used to validate key findings from scRNA-seq and CyTOF analyses.

\subsection{Computational analysis}
scRNA-seq data were processed using Cell Ranger, normalized with Scanpy, and clustered using the Leiden algorithm. TCR and BCR sequences were assembled and clonotypes defined using scRepertoire. Clonality was calculated as the normalized Shannon entropy. Cell--cell interactions were inferred using CellChat. The immune aging clock was built using elastic net regression on cell type proportions and transcriptomic features, with leave-one-out cross-validation.

\section{Results}

\subsection{Immune cell composition changes dynamically across the lifespan}
We profiled peripheral immune cells from 220 individuals spanning 13 age groups. Analysis of major immune cell populations revealed systematic age-associated changes (Table~\ref{tab:composition}). CD4$^+$ T cells declined progressively from 38.2\% in neonates (0--1 year) to 25.5\% in individuals over 80 years. In contrast, CD8$^+$ T cells increased from 12.5\% to 27.8\%, and NK cells expanded from 8.1\% to 22.5\% across the same range. B cells showed a marked decline from 22.8\% in neonates to 6.1\% in the oldest group. Monocytes remained relatively stable (14.3--15.5\%), while dendritic cells showed a modest decline from 4.1\% to 3.1\%.

\begin{table}[h]
\centering
\caption{Relative proportions (\%) of major immune populations by age group.}
\label{tab:composition}
\begin{tabular}{lcccccc}
\toprule
Age group & CD4$^+$ T & CD8$^+$ T & B cells & NK cells & Monocytes & DCs \\
\midrule
0--1 yr & 38.2 & 12.5 & 22.8 & 8.1 & 14.3 & 4.1 \\
1--5 yr & 36.5 & 14.1 & 20.3 & 9.5 & 15.2 & 4.4 \\
6--12 yr & 35.8 & 16.2 & 17.5 & 11.3 & 15.0 & 4.2 \\
13--18 yr & 34.2 & 18.6 & 14.8 & 12.8 & 15.3 & 4.3 \\
19--30 yr & 33.5 & 20.1 & 12.5 & 14.2 & 15.5 & 4.2 \\
31--45 yr & 32.8 & 21.5 & 11.2 & 15.5 & 15.0 & 4.0 \\
46--60 yr & 31.2 & 23.8 & 9.8 & 16.8 & 14.5 & 3.9 \\
61--70 yr & 29.5 & 25.2 & 8.5 & 18.2 & 15.0 & 3.6 \\
71--80 yr & 27.8 & 26.5 & 7.2 & 20.1 & 15.2 & 3.2 \\
$>$80 yr & 25.5 & 27.8 & 6.1 & 22.5 & 15.0 & 3.1 \\
\bottomrule
\end{tabular}
\end{table}

\subsection{T cell clonal expansion peaks in the elderly}
Analysis of paired TCR sequences revealed age-dependent clonal expansion patterns across T cell subsets (Table~\ref{tab:clonal}). GNLY$^+$CD8$^+$ effector memory T cells exhibited the highest clonality overall and the steepest age-associated increase, rising from a mean clonality of 0.08 in children to 0.52 in the elderly, with peak expansion in the $>$80 year group. CD4$^+$ central memory T cells and regulatory T cells also showed age-dependent clonal expansion, albeit less dramatically. MAIT cells displayed a unique trajectory, peaking in the 31--45 year group before declining in the elderly.

\begin{table}[h]
\centering
\caption{Mean clonality of T cell subsets across age categories.}
\label{tab:clonal}
\begin{tabular}{lcccc}
\toprule
T cell subset & Children & Adults & Elderly & Peak age \\
\midrule
GNLY$^+$CD8$^+$ EM T cells & 0.08 & 0.25 & 0.52 & $>$80 yr \\
CD8$^+$ naive T cells & 0.01 & 0.02 & 0.04 & 71--80 yr \\
CD4$^+$ central memory & 0.03 & 0.09 & 0.15 & 61--70 yr \\
Treg cells & 0.02 & 0.05 & 0.08 & $>$80 yr \\
MAIT cells & 0.12 & 0.18 & 0.14 & 31--45 yr \\
\bottomrule
\end{tabular}
\end{table}

\subsection{Cytotoxic B cells are enriched in children}
We identified a previously uncharacterized population of B cells expressing cytotoxic molecules, including granzyme B (GZMB). These cytotoxic B cells were markedly enriched in early life, constituting 5.2\% of B cells in neonates (0--1 year) and 4.8\% in young children (1--5 years), declining progressively to only 0.5\% in the oldest age group (Table~\ref{tab:cytob}). GZMB expression levels followed a parallel trajectory, from 3.8 log$_2$ CPM in neonates to 0.8 log$_2$ CPM in the elderly. This population may play a role in early immune defense and warrants further functional characterization.

\begin{table}[h]
\centering
\caption{Cytotoxic B cell fraction and GZMB expression by age group.}
\label{tab:cytob}
\begin{tabular}{lcc}
\toprule
Age group & Cytotoxic B cells (\%) & GZMB (log$_2$ CPM) \\
\midrule
0--1 yr & 5.2 & 3.8 \\
1--5 yr & 4.8 & 3.5 \\
6--12 yr & 3.1 & 2.9 \\
19--30 yr & 1.2 & 1.5 \\
46--60 yr & 0.8 & 1.1 \\
$>$80 yr & 0.5 & 0.8 \\
\bottomrule
\end{tabular}
\end{table}

\subsection{A single-cell immune aging clock}
We constructed a single-cell immune aging clock by integrating 156 features across 28 cell types, including cell type proportions, transcriptomic signatures, and clonal diversity metrics. Using elastic net regression with leave-one-out cross-validation, the clock achieved a Pearson correlation of $r = 0.92$ between predicted and chronological age and a mean absolute error of 5.8 years. This tool enables quantitative assessment of individual immune aging status and may identify individuals with accelerated or decelerated immune aging.

\section{Discussion}

Our multi-omic atlas of human peripheral immune cells across the full lifespan reveals several key findings. First, the immune system undergoes continuous, largely monotonic remodeling from birth to old age, with T cells being the most affected lineage. The progressive decline of CD4$^+$ T cells and B cells, coupled with expansion of CD8$^+$ T cells and NK cells, is consistent with previous reports but is here quantified at unprecedented resolution across the full age spectrum \citep{nikolich2018immune, goronzy2013understanding}.

The identification of GNLY$^+$CD8$^+$ effector memory T cells as the most clonally expanded subset in aging suggests that chronic antigen exposure or homeostatic proliferation drives the accumulation of large cytotoxic clones. The 6.5-fold increase in clonality from children to the elderly is striking and may contribute to the reduced T cell diversity observed in immunosenescence \citep{britanova2014age}.

The discovery of cytotoxic B cells enriched in early life is novel and suggests that the neonatal and pediatric immune system employs unconventional effector mechanisms. Granzyme B$^+$ B cells have been described in certain pathological contexts \citep{glass2020unconventional}, but their enrichment in healthy children has not been previously reported. These cells may contribute to early immune defense against pathogens or to immune regulation during immune system maturation.

The immune aging clock demonstrates that single-cell multi-omic data can predict chronological age with high accuracy. The mean absolute error of 5.8 years is comparable to or better than existing clocks based on DNA methylation or plasma proteomics \citep{sayed2021immuneclock}, and the single-cell resolution enables identification of the cell types and features that contribute most to the aging signature.

A limitation of this study is that all participants were from a single Chinese cohort, and the generalizability of findings to other populations requires validation. Additionally, the cross-sectional design cannot distinguish age effects from cohort effects.

\section{Conclusion}

We present a comprehensive multi-omic single-cell atlas of human peripheral immune aging spanning from birth to over 90 years of age. Our findings reveal that T cells are the most dynamically affected lineage, with GNLY$^+$CD8$^+$ effector memory T cells exhibiting the highest clonal expansion. We discover cytotoxic B cells enriched in early life and construct an immune aging clock with high predictive accuracy. This resource provides a foundation for future studies of immune aging, vaccine design, and age-related immune dysfunction.

\bibliographystyle{plainnat}
\bibliography{references}

\end{document}
